\section{NLP 模型工具箱的演进}

Manning 在此处回顾了 NLP 中可用的建模工具，形成了一个从简单到复杂的技术谱系：

\begin{figure}[H]
\centering
\includegraphics[width=0.95\textwidth]{slides-images/slide-026.jpg}
\caption{NLP 建模工具箱：从词袋模型到 Transformer 的演进\protect\footnotemark}
\end{figure}
\footnotetext{Slides 第27页。}

\begin{enumerate}
    \item \textbf{词向量 + 词袋模型}：最简单的分类方式，忽略词序
    \item \textbf{窗口模型}：类似 CNN 但更 ad hoc，早期课程中介绍
    \item \textbf{CNN}：擅长分类任务，容易并行化
    \item \textbf{RNN}：认知上更合理（从左到右阅读），但\textbf{难以并行化}
    \item \textbf{Transformer}：目前 NLP 的主流，Vision 领域也在越来越多地使用
\end{enumerate}

\begin{knowledgebox}{Batch Normalization vs Layer Normalization}
Manning 补充了一个关于归一化的知识点。在 CNN 中通常使用 \textbf{Batch Normalization}（先于 Layer Normalization 被发明），而 Transformer 中使用 \textbf{Layer Normalization}：
\begin{itemize}
    \item \textbf{Layer Norm}：在特征维度上计算统计量（对单个样本的所有特征做归一化）
    \item \textbf{Batch Norm}：在 batch 维度上计算统计量（对 batch 中所有样本的同一特征做归一化）
\end{itemize}
两者都实现了``零均值、单位方差''的标准化效果，但在不同架构中各有优势。
\end{knowledgebox}

\subsection{1x1 卷积的意义}

一个看似无意义的概念：\textbf{1x1 卷积}（size-1 convolution）。它实际上等价于对每个位置独立地施加一个全连接层，即：

$$\mathbf{h}_i = f(W\mathbf{x}_i + b)$$

这与 Transformer 中的 \textbf{Feed-Forward Layer}（FFN）完全相同——在每个 token 位置独立地进行非线性变换。1x1 卷积的参数量远少于跨越整个序列的全连接层。

\begin{importantbox}{1x1 卷积 = Position-wise FFN}
Transformer 中每个位置独立的前馈网络（Feed-Forward Network），从卷积的视角来看，就是一个 1x1 卷积。这种``在单个位置进行非线性变换''的操作是一个跨架构的通用设计模式。
\end{importantbox}

\subsection{本章小结}

NLP 模型工具箱经历了从词袋模型到 Transformer 的演进。不同架构各有优势：CNN 擅长分类和并行化，RNN 更符合人类阅读直觉，Transformer 兼具两者优点。Batch Norm 和 Layer Norm 是不同架构中常用的归一化技术，1x1 卷积与 Transformer 的 FFN 在数学上等价。

%% ========== VD-CNN ==========

